\subsection*{A.4. An admissible specialization}
For \cref{lem:admissible-nonempty} use
\[
 (c,\lambda,e_1,e_2,e_3)=(1,1,1,0,2),\qquad (u:v:z)=(3:5:1)
\]
over $\mathbf F_{13}$.  The shifted elliptic equation is
\[
 Y^2+2Y=x^3-6x^2+x+5,
\]
and the branch cubic $Y=0$ has discriminant $4$.  The fixed divisor is cut
out by
\[
 p(w)=w^3-w^2-2,
 \qquad \operatorname{disc}(p)=1,
\]
so it is reduced.  The residual conic has coefficient vector
\[
 (q_{00},q_{01},q_{02},q_{11},q_{12},q_{22})=(5,1,7,7,0,2)
\]
and factors as
\[
 5+x+7Y+7x^2+2Y^2
 =5(1+9x+Y)(1+12x+3Y).
\]
Thus the three norm lines are
\[
 L_0:Y=x,\qquad L_1:Y=4x-1,\qquad L_2:Y=4-4x.
\]
Substitution into the elliptic equation gives the three line-section cubics
\[
\begin{array}{c|c|c}
&\text{cubic in }x&\text{discriminant}\\ \hline
L_0&x^3+6x^2-x+5&4\\
L_1&x^3+4x^2+x+6&3\\
L_2&x^3+4x^2+2x-6&5.
\end{array}
\]
Hence all three intersections with $E$ are transverse.  The residual lines
avoid the branch divisor: $L_1\cap\{Y=0\}$ has $x=10$ and the elliptic
residual is $12\ne0$, while $L_2\cap\{Y=0\}$ has $x=1$ and residual
$1\ne0$.  Their pairwise intersections are
\[
 L_1\cap L_2=(12,8),\qquad
 L_0\cap L_1=(9,9),\qquad
 L_0\cap L_2=(6,6),
\]
and substitution in the elliptic equation gives residuals $5,11,11$,
respectively.  Thus no two of the three degree-three line sections meet on
$E$.

The two residual line sections are therefore finite \emph{\'etale} away from
the branch locus and occur with multiplicity one in the norm divisor.  The
local product formula for the norm then forces exactly one of the three
sheets of $C\to E$ to carry a simple zero above each residual point, which is
condition~(iii) of \cref{def:admissible-locus}.

The remaining canonical-model input can be checked from the same functions
used in \cref{prop:theta-resolution}.
Because $\mathbf F_{13}$ contains $\zeta_3$, the nine quadratic
representatives split into
\[
 \langle1,x,Y,x^2\rangle,\qquad
 \langle w,xw,wY\rangle,\qquad
 \langle w^2,xw^2\rangle.
\]
Their pole orders at a point over $O$ are $(0,2,3,4)$, $(1,3,4)$, and
$(2,4)$, respectively.  Thus the only canonical quadric is
$X_0X_2-X_1^2=0$, of rank three, and $|\kappa|$ is the unique $g^1_3$.
A member of this pencil is cut out by $a+bw$, so
$\Nm_\pi(a+bw)=a^3+b^3Y$; in particular its norm line has zero
$x$-coefficient.  By contrast, $L_0,L_1,L_2$ all have nonzero
$x$-coefficient.  Their inverse degree-three
divisors are therefore not linearly equivalent to $\kappa$, hence lie in the
$h^0=1$ Abel locus and map to the smooth theta locus.  This verifies
condition~(iv) and completes the arithmetic specialization used in
\cref{lem:admissible-nonempty}.

\begin{thebibliography}{99}

\bibitem{BBD1982}
A.~A. Beilinson, J. Bernstein, and P. Deligne,
\emph{Faisceaux pervers}, Ast\'erisque \textbf{100}, Soci\'et\'e Math\'ematique de France, Paris, 1982.

\bibitem{BirkenhakeLange2004}
Christina Birkenhake and Herbert Lange,
\emph{Complex abelian varieties}, second ed., Grundlehren der mathematischen Wissenschaften, vol.~302, Springer-Verlag, Berlin, 2004.

\bibitem{BrasseletLeSeade2000}
Jean-Paul Brasselet, D\~ung Tr\'ang L\^e, and Jos\'e Seade,
\emph{Euler obstruction and indices of vector fields}, Topology \textbf{39} (2000), no.~6, 1193--1208.

\bibitem{deCataldoMigliorini2002}
Mark Andrea A. de Cataldo and Luca Migliorini,
\emph{The hard Lefschetz theorem and the topology of semismall maps}, Ann. Sci. \`Ecole Norm. Sup. (4) \textbf{35} (2002), no.~5, 759--772.

\bibitem{Dynkin1957}
E.~B. Dynkin,
\emph{Semisimple subalgebras of semisimple Lie algebras}, Amer. Math. Soc. Transl. Ser.~2 \textbf{6} (1957), 111--244.

\bibitem{FraneckiKapranov2000}
J. Franecki and M. Kapranov,
\emph{The Gauss map and a noncompact Riemann--Roch formula for constructible sheaves on semiabelian varieties}, Duke Math. J. \textbf{104} (2000), no.~1, 171--180.

\bibitem{FultonHarris1991}
William Fulton and Joe Harris,
\emph{Representation theory: A first course}, Graduate Texts in Mathematics, vol.~129, Springer-Verlag, New York, 1991.

\bibitem{EGA4}
Alexander Grothendieck and Jean Dieudonn\'e,
\emph{\'El\'ements de g\'eom\'etrie alg\'ebrique. IV. \'Etude locale des sch\'emas et des morphismes de sch\'emas, troisi\`eme partie}, Publ. Math. Inst. Hautes \'Etudes Sci. \textbf{28} (1966), 5--255.

\bibitem{Hartshorne1977}
Robin Hartshorne,
\emph{Algebraic geometry}, Graduate Texts in Mathematics, vol.~52, Springer-Verlag, New York, 1977.

\bibitem{JKLM2025}
Ariyan Javanpeykar, Thomas Kr\"amer, Christian Lehn, and Marco Maculan,
\emph{The monodromy of families of subvarieties on abelian varieties}, Duke Math. J. \textbf{174} (2025), no.~6, 1045--1149.

\bibitem{Kashiwara1985}
Masaki Kashiwara,
\emph{Index theorem for constructible sheaves}, in \emph{Syst\`emes diff\'erentiels et singularit\'es}, Ast\'erisque \textbf{130} (1985), 193--209.

\bibitem{Kempf1973}
George Kempf,
\emph{On the geometry of a theorem of Riemann}, Ann.\ of Math. (2) \textbf{98} (1973), no.~1, 178--185.

\bibitem{KLM2026}
Thomas Kr\"amer, Christian Lehn, and Marco Maculan,
\emph{Exceptional Tannaka groups only arise from cubic threefolds}, J. Reine Angew. Math. \textbf{830} (2026), 51--99.

\bibitem{KraemerWeissauer2015}
Thomas Kr\"amer and Rainer Weissauer,
\emph{Vanishing theorems for constructible sheaves on abelian varieties}, J. Algebraic Geom. \textbf{24} (2015), no.~3, 531--568.

\bibitem{Laufer1977}
Henry B. Laufer,
\emph{On minimally elliptic singularities}, Amer. J. Math. \textbf{99} (1977), no.~6, 1257--1295.

\bibitem{Litt13}
Daniel Litt,
\emph{Problem no.~13}, Problems I Like, \url{https://www.problemsilike.com/13}, accessed August 10, 2026.

\bibitem{Macdonald1962}
I.~G. Macdonald,
\emph{Symmetric products of an algebraic curve}, Topology \textbf{1} (1962), 319--343.

\bibitem{Pinkham1977}
Henry C. Pinkham,
\emph{Simple elliptic singularities, del Pezzo surfaces and Cremona transformations}, Several Complex Variables, Proc. Sympos. Pure Math., vol.~30, American Mathematical Society, Providence, RI, 1977, pp.~69--71.

\bibitem{Saito1988}
Morihiko Saito,
\emph{Modules de Hodge polarisables}, Publ. Res. Inst. Math. Sci. \textbf{24} (1988), no.~6, 849--995.

\bibitem{stacks-project}
The Stacks Project Authors,
\emph{Stacks Project}, \url{https://stacks.math.columbia.edu}, accessed August 10, 2026.

\end{thebibliography}

\end{document}


